% This is samplepaper.tex, a sample chapter demonstrating the
% LLNCS macro package for Springer Computer Science proceedings;
% Version 2.21 of 2022/01/12
%
\documentclass[runningheads]{llncs}
%
\usepackage[T1]{fontenc}
\usepackage{graphicx}
\usepackage{hyperref}

\begin{document}
%
\title{Automated Pipeline for High-Fidelity 3D Visualization of Multi-Modal MRI Neuroimaging via Deep Learning and Real-Time Rendering Engines}
%
\titlerunning{Automated MRI-to-Unreal Pipeline}
%
\author{Dipenshu Deep Bhat\inst{1} \and
Anik Ghosh\inst{1}}
%
\authorrunning{F. Author et al.}
% First names are abbreviated in the running head.
% If there are more than two authors, 'et al.' is used.
%
\institute{Department of Biosciences and Bioengineering, Indian Institute of Technology Jammu, Jammu, India\\
\email{\{2025pmd00xx, 2025pmd0054\}@iitjammu.edu.in}}
%
\maketitle              % typeset the header of the contribution
%
\begin{abstract}
The transition from volumetric medical imaging data to interactive three-dimensional (3D) environments remains a significant bottleneck in surgical planning and medical education. This research presents an end-to-end automated pipeline that integrates deep learning-based segmentation with real-time rendering technologies. Utilizing the Medical Open Network for AI (MONAI) framework, we employ the BraTS and UNEST architectures to achieve high-precision segmentation of both pathological tumor regions and 132 distinct anatomical brain structures from multi-modal MRI scans. A novel processing bridge transforms segmented voxel data into refined polygonal meshes using the Marching Cubes algorithm, optimized for real-time performance through geometric smoothing. The pipeline culminates in a custom automation script for Unreal Engine, which utilizes a metadata-driven manifest to programmatically reconstruct the 3D scene. This system automates asset discovery, spatial positioning, and dynamic material assignment. The results demonstrate a significant reduction in visualization latency, offering a robust framework for immersive medical analysis and enhanced diagnostic interpretation.

\keywords{Medical Imaging \and MRI Segmentation \and Unreal Engine \and Deep Learning \and 3D Visualization \and MONAI.}
\end{abstract}
%
%
%
\section{Introduction}
Modern neuroimaging produces vast amounts of volumetric data, yet the interpretation of complex spatial relationships—such as tumor infiltration relative to critical anatomical structures—remains challenging when viewed on 2D screens. While 3D reconstruction exists, it often requires intensive manual labor or lacks the interactivity provided by modern game engines. This paper proposes a fully automated pipeline to convert NIfTI format MRI data into an interactive, color-coded 3D scene within Unreal Engine.

\section{Methodology}

\subsection{Deep Learning Segmentation}
The first phase utilizes the MONAI framework within a Jupyter environment. We implement two primary models:
\begin{itemize}
    \item \textbf{BraTS (Brain Tumor Segmentation):} Analyzes multi-modal sequences (FLAIR, T1, T2) to segment the necrotic core, edema, and enhancing tumor.
    \item \textbf{UNEST:} Conducts whole-brain segmentation to identify 132 anatomical regions from T1-weighted scans.
\end{itemize}

\subsection{Voxel-to-Mesh Conversion}
Post-segmentation, the pipeline employs the Marching Cubes algorithm to convert voxel-based masks into triangular meshes. To ensure engine compatibility, the \textit{trimesh} library is used for hole filling, normal correction, and surface smoothing. A metadata manifest (JSON) is generated to store medical labels and diagnostic color mappings.

\subsection{Unreal Engine Automation}
The final phase leverages the Unreal Engine Python API. The script (\texttt{unreal.py}) performs three critical functions:
\begin{enumerate}
    \item \textbf{Manifest Loading:} Reads JSON maps to identify required assets.
    \item \textbf{Procedural Spawning:} Instantiates StaticMeshActors at specific spatial coordinates to maintain anatomical integrity.
    \item \textbf{Dynamic Styling:} Creates Material Instances for each region, applying specific transparency to the brain shell and diagnostic colors to tumor segments.
\end{enumerate}

\section{Results and Discussion}
The automated pipeline successfully generated a 3D environment containing over 130 distinct anatomical regions and three tumor classifications in under ten minutes of processing time. By utilizing a "manifest" approach, we eliminated manual asset naming and placement errors. The resulting scene allows for real-time exploration, providing a spatial clarity that traditional slicing methods lack.

\section{Conclusion}
This research demonstrates that combining medical AI with real-time rendering engines can significantly streamline clinical visualization workflows. Future work will focus on integrating real-time haptic feedback for surgical simulation within the generated environment.

%
% ---- Bibliography ----
%
\begin{thebibliography}{8}
\bibitem{ref_monai}
Cardoso, M.J., et al.: MONAI: An open-source framework for deep learning in healthcare. arXiv preprint arXiv:2211.02701 (2022).

\bibitem{ref_brats}
Menze, B.H., et al.: The Multimodal Brain Tumor Image Segmentation Benchmark (BRATS). IEEE Transactions on Medical Imaging 34(10), 1993--2024 (2015).

\bibitem{ref_unreal}
Unreal Engine Python API Documentation. Epic Games (2024).

\end{thebibliography}
\end{document}